\documentclass[aspectratio=169,10pt]{beamer}
\usetheme[numbering=fraction,progressbar=frametitle]{metropolis}
\usepackage{fontspec}
\setsansfont{Latin Modern Sans}
\usepackage{amsmath,amssymb,mathtools,graphicx,booktabs}
\usepackage{polyglossia}
\setdefaultlanguage{english}
\setbeamertemplate{navigation symbols}{}
\setbeamerfont{frametitle}{size=\large}
\setbeamerfont{normal text}{size=\normalsize}
\setlength{\parskip}{.45em}
\hypersetup{pdfauthor={},pdfsubject={Lecture slides},pdfkeywords={},pdftitle={Differentiation, convexity and continuity}}
\providecommand{\N}{\mathbb{N}}
\providecommand{\Z}{\mathbb{Z}}
\providecommand{\D}{\mathbb{D}}
\providecommand{\Q}{\mathbb{Q}}
\providecommand{\R}{\mathbb{R}}
\providecommand{\C}{\mathbb{C}}
\providecommand{\vect}[1]{\overrightarrow{#1}}
\providecommand{\vi}{\vec{\imath}}
\providecommand{\vj}{\vec{\jmath}}
\providecommand{\vk}{\vec{k}}
\providecommand{\norm}[1]{\left\lVert #1\right\rVert}
\providecommand{\vc}[2]{\begin{pmatrix}#1\\#2\end{pmatrix}}
\providecommand{\deter}[2]{\det\left(#1,#2\right)}
\providecommand{\abs}[1]{\left\lvert #1\right\rvert}
\providecommand{\dd}{\,\mathrm{d}}
\providecommand{\e}{\mathrm{e}}
\providecommand{\Card}{\operatorname{Card}}
\providecommand{\E}{\operatorname{E}}
\providecommand{\V}{\operatorname{V}}
\providecommand{\degre}{^{\circ}}
\providecommand{\un}[1]{\,\mathrm{#1}}
\providecommand{\sci}[2]{#1\times10^{#2}}
\providecommand{\pgcd}{\operatorname{PGCD}}
\providecommand{\ppcm}{\operatorname{PPCM}}
\title{Differentiation, convexity and continuity}
\subtitle{Lecture and worked examples}
\date{}
\author{}
\begin{document}
\begin{frame}[plain]
\vfill
{\usebeamerfont{title}\inserttitle\par}
\medskip
{\usebeamerfont{subtitle}\insertsubtitle\par}
\medskip
{\usebeamercolor[fg]{progress bar}\rule{\textwidth}{.5pt}}
\vfill
\end{frame}
\begin{frame}{Objectives}
Differentiate a composite function; study variations and convexity; prove existence and uniqueness; approximate a solution.
\end{frame}
\begin{frame}{Differentiating a composite}
If \(u\) is differentiable on \(I\) and \(v\) is differentiable on an interval containing \(u(I)\), then \((v\circ u)'=(v'\circ u)u'\). Domains must be checked before computing.
\[ (u^n)'=nu'u^{n-1},\quad(e^u)'=u'e^u,\quad
(\dfrac{1}{u})'=-\dfrac{u'}{u^2},\quad(\sqrt u)'=\dfrac{u'}{2\sqrt u}\quad(u>0).\]
For example, \(f(x)=\sqrt{3x+1}\) is differentiable for \(x>-\dfrac{1}{3}\) and \(f'(x)=\dfrac{3}{2\sqrt{3x+1}}\). The sign of \(f'\) determines the variations of \(f\); the sign of \(f\) does not.
\end{frame}
\begin{frame}{Convexity}
A function \(f\) is convex on an interval \(I\) if its curve lies below its chords: for \(x,y\in I\) and \(t\in[0,1]\),
\[f((1-t)x+ty)\le(1-t)f(x)+tf(y).\]
For a twice-differentiable function, the following equivalences are accepted without proof: \(f\) is convex; \(f'\) is increasing; \(f''\ge0\); the curve lies above all its tangents. The inequalities are reversed for a concave function. A point of inflection is a point where the curve changes convexity.
\end{frame}
\begin{frame}{Convexity}
Suppose \(f''\ge0\). Fix \(a\) and set \(g(x)=f(x)-f(a)-f'(a)(x-a)\). Since \(f'\) is increasing, \(g'(x)=f'(x)-f'(a)\le0\) for \(x\le a\) and \(\ge0\) for \(x\ge a\). Thus \(g\) attains its minimum \(g(a)=0\) at \(a\), so \(f(x)\ge f(a)+f'(a)(x-a)\).
\end{frame}
\begin{frame}{Convexity}
\(x^4\) is convex and \(f''(0)=0\), but \(0\) is not the abscissa of a point of inflection. \(x^3\) changes convexity at \(0\). The convexity of the exponential and its tangent at \(0\) give \(e^x\ge1+x\).
\end{frame}
\begin{frame}{Continuity and equations}
\(f\) is continuous at \(a\) if \(\lim_{x\to a}f(x)=f(a)\). Every differentiable function is continuous. The converse is false: \(|x|\) is continuous at \(0\) but not differentiable at \(0\).
\end{frame}
\begin{frame}{Continuity and equations}
If \(f\) is continuous on \([a,b]\), every real number between \(f(a)\) and \(f(b)\) has at least one preimage in \([a,b]\). If \(f\) is also strictly monotonic, this preimage is unique. One can also work on an open interval using the limits at the endpoints.
\end{frame}
\begin{frame}{Continuity and equations}
\(f(x)=x^3+x-1\) is continuous on \([0,1]\). Since \(f(0)=-1\) and \(f(1)=1\), the IVT gives a zero in \(]0,1[\). Because \(f'(x)=3x^2+1>0\), it is unique. \(f(0.68)<0<f(0.69)\) provides a bracketing interval, without replacing the justifications of existence and uniqueness.
\end{frame}
\begin{frame}{Continuity and equations}
For \(x\ne1\), set \(g(x)=\dfrac{x^2-1}{x-1}\), and set \(g(1)=a\). Since \(g(x)=x+1\) when \(x\ne1\), \(\lim_{x\to1}g(x)=2\). The function is continuous at \(1\) if and only if \(a=2\). The value at the point must therefore be compared with the limit, not replaced by it without justification.
\end{frame}
\begin{frame}{Continuity and equations}
If \(u_n\to\ell\) and \(f\) is continuous at \(\ell\), then \(f(u_n)\to f(\ell)\). For a convergent recursive sequence \(u_{n+1}=f(u_n)\), we obtain \(\ell=f(\ell)\).
\end{frame}
\begin{frame}{Bisection}
\begin{columns}[T,onlytextwidth]
\begin{column}{.57\textwidth}
For \(f\) continuous and \(f(a)f(b)<0\), compute \(m=\dfrac{a+b}{2}\). If \(f(m)=0\), stop; otherwise keep the half whose endpoints have images of opposite signs. After \(n\) steps, the length is \(\dfrac{b-a}{2^n}\). Taking the midpoint gives an error of at most half the length. The method does not guarantee uniqueness; this is proved separately.   \(y=x^3-3x\) and its tangent at the point of inflection.
\end{column}
\begin{column}{.4\textwidth}
\centering
\includegraphics[width=\linewidth,height=4.5cm,keepaspectratio]{figures/terminale-06-1-en.pdf}
\end{column}
\end{columns}
\end{frame}
\begin{frame}{Research activity}
\textbf{Derivatives and graphical readings - 35 min.}

Consider the function \(f\) defined on \(\R\) by \(f(x)=x^3-3x\). Let \(\mathcal C_f\) denote its graph.
\end{frame}
\begin{frame}{Research activity (2)}
\begin{enumerate}
\setcounter{enumi}{0}\item Calculate \(f' (x)\) and \(f'' (x)\), then draw up their sign tables.

\item Deduce the variations of \(f\) and its intervals of convexity and concavity.

\item Find the coordinates of the point of inflection of \(\mathcal C_f\) and an equation of the tangent at that point.

\item Using limits, extrema and the intermediate value theorem, determine the number of real solutions of \(f(x)=1\).

\item Explain, using this example, why continuity is not enough to guarantee the uniqueness of a solution.

\end{enumerate}
\end{frame}
\end{document}
